% Preamble template for Tufte-style applied econometrics lecture notes.
% Do not compile this file directly; compile the main notes file instead.
\documentclass{tufte-handout}

\usepackage[T1]{fontenc}
\usepackage[utf8]{inputenc}
\usepackage{ntheorem}
\usepackage{graphicx}
\usepackage{amsmath}
\usepackage{amssymb}
\usepackage{mathtools}
\usepackage{hyperref}
\usepackage{epigraph}
\usepackage{booktabs}
\usepackage{array}
\theoremstyle{break}

\hypersetup{
  colorlinks,
  urlcolor = blue,
  pdfauthor={Swapnil Singh},
  pdfkeywords={econometrics, causal inference, applied econometrics},
  pdftitle={Lecture Notes for Applied Econometrics},
  pdfpagemode=UseNone
}

\newtheorem{ruleN}{Rule}
\newtheorem{thmN}{Theorem}
\newtheorem{assN}{Assumption}
\newtheorem{defN}{Definition}
\newtheorem{exmp}{Example}
\newtheorem{cmt}{Comment}
\newtheorem{proof}{Proof}

\newcommand{\continuation}{??}
\newtheorem*{excont}{Example \continuation}
\newenvironment{continueexample}[1]
 {\renewcommand{\continuation}{\ref{#1}}\excont[continued]}
 {\endexcont}

\newcommand{\bY}{\mathbf{Y}}
\newcommand{\bX}{\mathbf{X}}
\newcommand{\bD}{\mathbf{D}}
\newcommand{\E}{\mathbb{E}}
\newcommand{\Prob}{\mathbb{P}}
\newcommand{\Var}{\mathrm{Var}}
\newcommand{\Cov}{\mathrm{Cov}}
\newcommand{\se}{\mathrm{SE}}
\newcommand{\CI}{\mathrm{CI}}
\newcommand{\ATE}{\mathrm{ATE}}
\newcommand{\ATT}{\mathrm{ATT}}
\newcommand{\TOT}{\mathrm{TOT}}
\newcommand{\LATE}{\mathrm{LATE}}
\newcommand{\ITT}{\mathrm{ITT}}
\newcommand{\RD}{\mathrm{RD}}
\newcommand{\DiD}{\mathrm{DiD}}
\newcommand{\MMSE}{\mathrm{MMSE}}
\newcommand{\OLS}{\mathrm{OLS}}
\newcommand{\Normal}{\mathcal{N}}
\newcommand{\plim}{\operatorname*{plim}}
\DeclareMathOperator*{\argmin}{arg\,min}
\DeclareMathOperator{\Supp}{Supp}

\newcommand\independent{\protect\mathpalette{\protect\independenT}{\perp}}
\def\independenT#1#2{\mathrel{\rlap{$#1#2$}\mkern2mu{#1#2}}}
\newcommand{\indep}{\perp\!\!\!\perp}

\usepackage{cleveref}
\crefname{appsec}{appendix}{appendices}
\crefname{appsubsec}{appendix}{appendices}
\crefname{assumption}{assumption}{assumptions}
\crefname{equation}{equation}{equations}
\crefname{exmp}{example}{examples}
\crefname{assN}{assumption}{assumptions}
\crefname{cmt}{comment}{comments}
\crefname{defN}{definition}{definitions}

\usepackage[nolist]{acronym}
\begin{acronym}
  \acro{ATE}{average treatment effect}
  \acro{ATT}{average treatment effect on the treated}
  \acro{CIA}{conditional independence assumption}
  \acro{CI}{confidence interval}
  \acro{CLT}{central limit theorem}
  \acro{DiD}{difference-in-differences}
  \acro{FE}{fixed effects}
  \acro{FWL}{Frisch-Waugh-Lovell}
  \acro{IV}{instrumental variables}
  \acro{LATE}{local average treatment effect}
  \acro{OLS}{ordinary least squares}
  \acro{OVB}{omitted-variable bias}
  \acro{PO}{potential outcomes}
  \acro{RCT}{randomized controlled trial}
  \acro{RD}{regression discontinuity}
  \acro{SUTVA}{stable unit treatment value assignment}
  \acro{TWFE}{two-way fixed effects}
\end{acronym}

\def\inprobHIGH{\,{\buildrel p \over \rightarrow}\,}
\def\inprob{\,{\inprobHIGH}\,}
\def\indistHIGH{\,{\buildrel d \over \rightarrow}\,}
\def\indist{\,{\indistHIGH}\,}

\usepackage[many]{tcolorbox}
\definecolor{main}{HTML}{3F6EA7}
\definecolor{sub}{HTML}{F2F7FD}
\definecolor{sub2}{HTML}{FFF5E8}
\definecolor{mainline}{HTML}{D3DFEE}
\definecolor{warn}{HTML}{C7771B}
\definecolor{warnline}{HTML}{EAD9C0}

\tcbset{
    enhanced,
    breakable,
    colback = white,
    boxrule = 0.5pt,
    arc = 2pt,
    left = 9pt,
    right = 9pt,
    top = 7pt,
    bottom = 7pt,
    before skip = 0.3cm,
    after skip = 0.45cm
}

\newtcolorbox{boxD}{
    colback = sub,
    colframe = mainline,
    borderline west = {2.2pt}{0pt}{main},
    borderline north = {0.6pt}{0pt}{main!25},
    borderline south = {0.6pt}{0pt}{main!15}
}

\newtcolorbox{boxF}{
    colback = sub2,
    colframe = warnline,
    borderline west = {2.2pt}{0pt}{warn},
    borderline north = {0.6pt}{0pt}{warn!25},
    borderline south = {0.6pt}{0pt}{warn!15}
}

\usepackage{colortbl}
\usepackage{tikz}
\usepackage{verbatim}
\usetikzlibrary{positioning}
\usetikzlibrary{snakes}
\usetikzlibrary{calc}
\usetikzlibrary{arrows}
\usetikzlibrary{decorations.markings}
\usetikzlibrary{shapes.misc}
\usetikzlibrary{matrix,shapes,arrows,fit,tikzmark}


\title{Applied Econometrics: Session 5}
\author{Swapnil Singh}
\date{}

\hypersetup{
  pdftitle={Applied Econometrics: Session 5},
  pdfauthor={Swapnil Singh},
  pdfkeywords={instrumental variables, non-compliance, LATE, 2SLS, weak instruments}
}

\setlength{\headheight}{26pt}

\newcommand{\Zi}{Z_i}
\newcommand{\Di}{D_i}
\newcommand{\Yi}{Y_i}
\newcommand{\Xcov}{X_i}
\newcommand{\Ui}{U_i}
\newcommand{\Vi}{V_i}
\newcommand{\Yone}{Y_i(1)}
\newcommand{\Yzero}{Y_i(0)}
\newcommand{\Done}{D_i(1)}
\newcommand{\Dzero}{D_i(0)}
\newcommand{\Covp}{\operatorname{Cov}}
\newcommand{\Corr}{\operatorname{Corr}}

\begin{document}
\maketitle

\begin{abstract}
This lecture studies instrumental variables, non-compliance, and the local
average treatment effect. The central problem is hidden confounding: treatment
choice is related to unobserved determinants of the outcome. An instrument gives
us a source of variation in treatment that is as good as randomly assigned, but
only after we accept strong assumptions. We derive the first stage, reduced
form, Wald estimator, two-stage least squares, the weak-instrument problem, and
the interpretation of IV estimates in randomized experiments with imperfect
take-up.
\end{abstract}

\begin{boxD}
\textbf{Reading and objective.} Read Chapters 08 and 09 of the Python Causality
Handbook: Instrumental Variables and Non-Compliance/LATE. The goal is to know
when IV is useful, what assumptions make it credible, how to compute the
estimand, and how to explain the result in plain language. The companion Python
script is \texttt{lectures/code/lecture-5-iv-late.py}; run it from the course
root with \texttt{python lectures/code/lecture-5-iv-late.py}.
\end{boxD}

\section{Why Instrumental Variables?}

The basic causal question is familiar. We observe an outcome $\Yi$, a treatment
$\Di$, and perhaps controls $\Xcov$. We would like to estimate the effect of
treatment on the outcome. If we write the structural equation
\[
    \Yi = \alpha + \tau \Di + \Ui,
\]
then $\tau$ is the causal effect only if the treatment is unrelated to all
unobserved determinants of the outcome:
\[
    \Covp(\Di,\Ui)=0.
\]
This condition is often too strong. People choose treatments. Schools choose
programs. Firms choose policies. Governments target interventions. These choices
usually depend on information that the econometrician does not observe.

\begin{exmp}[Education and earnings]
Suppose $\Di$ is years of schooling and $\Yi$ is log earnings. Ability,
parental support, motivation, health, and local labor-market opportunities may
affect both schooling and earnings. If these variables are omitted, an OLS
regression of earnings on schooling mixes the causal return to schooling with
selection.
\end{exmp}

\subsection{Hidden confounding}

Hidden confounding means that the treated and untreated groups would have had
different outcomes even if nobody had received treatment. In potential-outcome
notation, the problem is
\[
    \E[\Yzero \mid \Di=1] \neq \E[\Yzero \mid \Di=0].
\]
The untreated outcome for the treated is missing. The untreated group is not a
valid replacement if selection into treatment is related to unobservables.

\begin{boxF}
\textbf{Common mistake.} Students often say ``we controlled for many things, so
there is no omitted-variable bias.'' That is not a theorem. Controls help only
for variables that are measured well and enter the model in a sufficiently rich
way. Unobserved ability remains unobserved even after a long control list.
\end{boxF}

\subsection{Instrument versus control}

An instrumental variable is not just another control. A control is a variable
we condition on because it is a confounder or helps make treatment as good as
random. An instrument is a variable that shifts treatment, but has no direct
path to the outcome except through treatment.

\begin{defN}[Control]
A control variable $\Xcov$ is included to compare treated and untreated units that
are similar in observed characteristics. The identifying hope is often
\[
    (\Yone,\Yzero) \independent \Di \mid \Xcov.
\]
\end{defN}

\begin{defN}[Instrument]
An instrument $\Zi$ is a variable that creates exogenous variation in treatment:
it affects $\Di$, is unrelated to unobserved outcome determinants, and affects
$\Yi$ only through $\Di$.
\end{defN}

The distinction is important. If distance to college affects schooling, it may
be an instrument for schooling. But if distance also affects earnings through
local labor-market quality, family background, or migration opportunities, it is
not a valid instrument. It may be a control, a proxy, or a bad variable, but it
does not automatically solve endogeneity.

\section{Valid Instruments}

The canonical linear IV model is
\begin{align}
    \Yi &= \alpha + \tau \Di + \Ui, \label{eq:structural}\\
    \Di &= \pi_0 + \pi_1 \Zi + \Vi. \label{eq:firststage}
\end{align}
Equation \eqref{eq:structural} is the causal equation we care about. Equation
\eqref{eq:firststage} is the first-stage equation: it tells us whether the
instrument moves treatment.

\begin{assN}[Relevance]
The instrument must predict treatment:
\[
    \Covp(\Zi,\Di) \neq 0.
\]
In the first-stage regression, this means $\pi_1 \neq 0$.
\end{assN}

Relevance is the easiest IV assumption to understand and the easiest to check.
If the instrument barely changes treatment, then the instrument gives little
useful causal variation. The denominator of the IV ratio becomes small, and the
estimate becomes unstable.

\begin{assN}[Independence]
The instrument must be as good as randomly assigned with respect to potential
outcomes and potential treatments. A common statement is
\[
    \Zi \independent (\Yone,\Yzero,\Done,\Dzero).
\]
In a linear structural equation, the corresponding moment is
\[
    \E[\Zi \Ui] = 0,
\]
or, after centering, $\Covp(\Zi,\Ui)=0$.
\end{assN}

Independence is not the same as relevance. A variable can be strongly related
to treatment and still be invalid. For example, parental education may strongly
predict a child's schooling, but it likely also predicts earnings through
family networks, wealth, and preferences.

\begin{assN}[Exclusion restriction]
The instrument affects the outcome only through the treatment. In potential
outcome language, changing $Z$ while holding treatment received fixed cannot
change $Y$.
\end{assN}

The exclusion restriction is the most fragile assumption. It is also usually
untestable. A lottery offer for a training program may satisfy exclusion if the
offer only matters by changing actual training. It fails if receiving the offer
also raises confidence, improves access to counselors, changes job-search
behavior, or signals quality to employers even for people who never attend.

\subsection{What can be tested?}

The first stage can be estimated. Balance on observed covariates can be checked
when the instrument is randomized or plausibly random. Overidentification tests
can reject some combinations of assumptions when there are more instruments than
endogenous variables. But the core exclusion restriction cannot be proven from
the same data used for the outcome regression.

\begin{boxD}
\textbf{Practical rule.} A credible IV design needs a story before it needs a
regression table. The story must explain why $Z$ moves $D$, why $Z$ is not
chosen in response to unobserved determinants of $Y$, and why $Z$ has no direct
effect on $Y$.
\end{boxD}

\subsection{Common failure modes}

Instruments fail in predictable ways.
\begin{itemize}
  \item The instrument is a proxy for omitted background characteristics.
  \item The instrument changes expectations, information, prices, peer groups,
        or effort directly, not only treatment.
  \item The instrument affects several treatments at once.
  \item The instrument is chosen by agents who anticipate potential outcomes.
  \item The instrument is weak, so finite-sample IV behaves badly.
\end{itemize}

\section{Wald IV}

Start with a binary instrument, $\Zi \in \{0,1\}$, and one treatment $\Di$.
The first stage is the effect of the instrument on treatment:
\[
    \text{First stage}
    =
    \E[\Di \mid \Zi=1] - \E[\Di \mid \Zi=0].
\]
The reduced form is the effect of the instrument on the outcome:
\[
    \text{Reduced form}
    =
    \E[\Yi \mid \Zi=1] - \E[\Yi \mid \Zi=0].
\]
The Wald estimator divides the reduced form by the first stage:
\[
    \hat{\tau}_{Wald}
    =
    \frac{\bar{Y}_{Z=1}-\bar{Y}_{Z=0}}
         {\bar{D}_{Z=1}-\bar{D}_{Z=0}}.
\]

\begin{thmN}[Wald IV in the constant-effect linear model]
Suppose
\[
    \Yi = \alpha + \tau \Di + \Ui,
    \qquad \Covp(\Zi,\Ui)=0,
    \qquad \Covp(\Zi,\Di)\neq 0.
\]
Then the population IV estimand is
\[
    \tau
    =
    \frac{\Covp(\Zi,\Yi)}{\Covp(\Zi,\Di)}.
\]
If $Z$ is binary, this equals
\[
    \tau
    =
    \frac{\E[\Yi\mid \Zi=1]-\E[\Yi\mid \Zi=0]}
         {\E[\Di\mid \Zi=1]-\E[\Di\mid \Zi=0]}.
\]
\end{thmN}

\begin{proof}
Take covariance of the structural equation with $Z$:
\[
    \Covp(\Zi,\Yi)
    =
    \Covp(\Zi,\alpha + \tau \Di + \Ui)
    =
    \tau\Covp(\Zi,\Di) + \Covp(\Zi,\Ui).
\]
By instrument exogeneity, $\Covp(\Zi,\Ui)=0$. If relevance holds, divide by
$\Covp(\Zi,\Di)$.
\end{proof}

\begin{exmp}[Numerical Wald calculation]
Suppose a scholarship offer $Z$ raises college attendance $D$ from $0.40$ to
$0.70$. The first stage is $0.30$. Suppose average earnings are $40$ for those
not offered and $43$ for those offered. The reduced form is $3$. The Wald IV
estimate is
\[
    \frac{43-40}{0.70-0.40}=10.
\]
The interpretation is: for each additional person induced into college by the
scholarship offer, earnings rise by $10$ units on average, under the IV
assumptions.
\end{exmp}

\subsection{Why the ratio makes sense}

The reduced form is an effect of the instrument, not an effect of treatment.
If a scholarship offer raises attendance by $30$ percentage points and earnings
by $3$ units, then the $3$-unit earnings effect is diluted across everyone
offered the scholarship, including people whose attendance did not change. IV
rescales the diluted outcome effect by the treatment take-up effect.

\section{IV with Controls}

Controls can be used in IV, but they do not rescue a bad instrument. With
controls, the assumptions become conditional:
\[
    \E[\Zi \Ui \mid \Xcov]=0,
    \qquad
    \Covp(\Zi,\Di \mid \Xcov) \neq 0.
\]
The structural and first-stage equations are
\begin{align}
    \Yi &= \alpha + \tau \Di + \Xcov'\beta + \Ui,\\
    \Di &= \pi_0 + \pi_1 \Zi + \Xcov'\gamma + \Vi.
\end{align}

One way to understand controls in IV is residualization. First remove from
$Y$, $D$, and $Z$ the part predicted by $X$. Then run IV using the residualized
variables. This is the Frisch-Waugh-Lovell logic applied to IV.

\begin{boxF}
\textbf{Bad control warning.} Do not control for variables caused by the
instrument or the treatment unless the research design explicitly calls for it.
Post-treatment controls can block part of the causal effect or open new
selection paths.
\end{boxF}

\section{Two-Stage Least Squares}

The Wald estimator is the simplest IV estimator. Two-stage least squares is the
standard regression implementation. It extends naturally to continuous
instruments, controls, and multiple instruments.

\subsection{Stage 1}

Regress treatment on the instrument and controls:
\[
    \Di = \pi_0 + \pi_1 \Zi + \Xcov'\gamma + \Vi.
\]
Compute the fitted values
\[
    \hat{D}_i = \hat{\pi}_0 + \hat{\pi}_1 \Zi + \Xcov'\hat{\gamma}.
\]
The fitted treatment $\hat{D}_i$ is the part of treatment predicted by the
instrument and controls. If the instrument is valid, this predicted part is the
exogenous part of treatment variation.

\subsection{Stage 2}

Regress the outcome on the fitted treatment and controls:
\[
    \Yi = \alpha + \tau \hat{D}_i + \Xcov'\beta + \text{residual}.
\]
The coefficient on $\hat{D}_i$ is the 2SLS estimate.

\begin{boxF}
\textbf{Inference warning.} It is fine to explain 2SLS as two regressions, but
do not manually run the two regressions and use the second-stage OLS standard
error. The generated regressor $\hat{D}$ makes those standard errors wrong.
Use an IV/2SLS routine that computes the correct variance.
\end{boxF}

\subsection{Projection formula}

Let $W=[\mathbf{1},D,X]$ be the matrix of structural regressors and let
$Q=[\mathbf{1},Z,X]$ be the matrix of instruments and included exogenous
variables. The 2SLS estimator is
\[
    \hat{\theta}_{2SLS}
    =
    (W'P_Q W)^{-1} W'P_Q Y,
    \qquad
    P_Q = Q(Q'Q)^{-1}Q'.
\]
The projection matrix $P_Q$ keeps only the part of the structural regressors
that is explained by the instruments and exogenous controls.

\subsection{Multiple instruments}

If there is one endogenous variable and several excluded instruments, 2SLS uses
all instruments to construct the predicted treatment. This can improve
precision if the instruments are valid and relevant. But extra instruments are
not free. Each additional instrument adds another exclusion restriction. A weak
or questionable extra instrument can make the design less credible.

\section{Weak Instruments}

An instrument is weak when it predicts treatment only weakly. Weak instruments
create two related problems. First, the estimator is noisy because the IV
denominator is small. Second, the finite-sample distribution of IV can be badly
non-normal and biased toward OLS.

For the just-identified Wald estimator,
\[
    \hat{\tau}_{IV}
    =
    \frac{\widehat{\text{reduced form}}}
         {\widehat{\text{first stage}}}.
\]
If the estimated first stage is close to zero, small sampling errors in the
denominator produce huge changes in the ratio.

\begin{defN}[First-stage F-statistic]
With one endogenous variable, the first-stage $F$-statistic tests whether the
excluded instruments jointly predict treatment after controls. In the simplest
case with one excluded instrument, it is the square of the first-stage
$t$-statistic on $Z$.
\end{defN}

The common rule of thumb is that a first-stage $F$ below about $10$ is a warning
sign. This is not a law of nature; it is a diagnostic. Serious empirical work
should report the first stage, discuss instrument strength, and consider
weak-IV robust inference when instrument strength is questionable.

\begin{boxD}
\textbf{Python example.} The script
\texttt{lectures/code/lecture-5-iv-late.py} simulates the same true treatment
effect with different first-stage strengths. As the first stage gets weaker,
the distribution of IV estimates becomes much wider even though the data
generating process has not changed.
\end{boxD}

\subsection{What IV does not fix}

IV is not a magic correction for every endogeneity problem. It does not fix a
direct effect of $Z$ on $Y$. It does not fix a non-random instrument. It does
not estimate the average treatment effect when treatment effects are
heterogeneous unless additional assumptions hold. It does not make a poor
measurement strategy credible. It answers the causal question created by the
instrument.

\section{Non-Compliance}

Randomized experiments often randomize assignment, not treatment received. Let
$Z_i$ denote assignment or offer, and let $D_i$ denote actual treatment
received. Non-compliance means
\[
    D_i \neq Z_i
\]
for some units. Some people assigned to treatment do not take it. Some people
assigned to control get access elsewhere.

\begin{exmp}[Job training offer]
A government randomly offers a training program. Some offered people do not
attend because of time, distance, or lack of interest. Some people not offered
find similar training elsewhere. Assignment is randomized, but attendance is a
choice.
\end{exmp}

\subsection{Why not compare by treatment received?}

Comparing treated and untreated people by actual receipt reintroduces selection.
People who comply with a training offer may be more motivated. People who seek
training without an offer may have better information. People who refuse
treatment may face constraints related to outcomes. Randomization protects $Z$,
not necessarily $D$.

\subsection{Intention-to-treat}

The intention-to-treat effect is the causal effect of assignment:
\[
    \ITT_Y
    =
    \E[\Yi \mid \Zi=1] - \E[\Yi \mid \Zi=0].
\]
It is often policy relevant. If a government can offer a program but cannot
force participation, the ITT tells us the effect of the offer policy.

There is also an ITT for treatment take-up:
\[
    \ITT_D
    =
    \E[\Di \mid \Zi=1] - \E[\Di \mid \Zi=0].
\]
This is the first stage in the non-compliance setting.

\section{Compliance Types}

Potential treatment receipt describes what each person would do under each
assignment:
\[
    \Dzero = \text{treatment received if } Z=0,
    \qquad
    \Done = \text{treatment received if } Z=1.
\]
With binary assignment and binary treatment, each unit belongs to one of four
latent compliance types.

\begin{center}
\begin{tabular}{lll}
\toprule
Type & $(\Dzero,\Done)$ & Interpretation \\
\midrule
Never-taker & $(0,0)$ & never takes treatment \\
Complier & $(0,1)$ & takes treatment only if assigned/offered \\
Always-taker & $(1,1)$ & takes treatment regardless of assignment \\
Defier & $(1,0)$ & does opposite of assignment \\
\bottomrule
\end{tabular}
\end{center}

These types are not usually observed at the individual level because each
person is seen under only one assignment. We observe either $D_i(0)$ or
$D_i(1)$, never both.

\begin{assN}[Monotonicity]
There are no defiers:
\[
    \Done \geq \Dzero
    \quad \text{for every } i.
\]
Assignment weakly increases the probability of treatment for everyone.
\end{assN}

Monotonicity is a behavioral assumption. It is plausible when an offer, subsidy,
or eligibility rule can only make access easier. It is less plausible when the
instrument changes incentives in different directions for different people.

\subsection{Which types move the first stage?}

Never-takers have $\Done-\Dzero=0$. Always-takers also have
$\Done-\Dzero=0$. Defiers have $\Done-\Dzero=-1$. Compliers have
$\Done-\Dzero=1$. Under monotonicity, the first stage equals the share of
compliers:
\[
    \E[\Di\mid Z_i=1] - \E[\Di\mid Z_i=0]
    =
    \Pr(\text{complier}).
\]
This is why IV with non-compliance estimates an effect for compliers. They are
the only group whose treatment status is changed by the instrument.

\section{Local Average Treatment Effect}

The local average treatment effect is the average treatment effect for
compliers:
\[
    \LATE
    =
    \E[\Yone-\Yzero \mid \Done > \Dzero].
\]
The word ``local'' does not mean local in geography. It means local to the
subpopulation whose treatment is changed by the instrument.

\begin{thmN}[LATE theorem]
Assume random assignment of $Z$, exclusion, relevance, and monotonicity. Then
the Wald estimand identifies the average treatment effect for compliers:
\[
    \frac{\E[\Yi\mid \Zi=1]-\E[\Yi\mid \Zi=0]}
         {\E[\Di\mid \Zi=1]-\E[\Di\mid \Zi=0]}
    =
    \E[\Yone-\Yzero \mid \Done>\Dzero].
\]
\end{thmN}

\begin{proof}
By exclusion, observed outcomes can be written as
\[
    Y_i = Y_i(0) + D_i\{Y_i(1)-Y_i(0)\}.
\]
Random assignment implies that differences by $Z$ reveal average differences
between potential outcomes under assignment. Therefore
\[
    \E[Y_i\mid Z_i=1]-\E[Y_i\mid Z_i=0]
    =
    \E[(D_i(1)-D_i(0))\{Y_i(1)-Y_i(0)\}].
\]
Under monotonicity, $D_i(1)-D_i(0)$ is one for compliers and zero for
never-takers and always-takers. Hence the numerator equals
\[
    \Pr(\text{complier})
    \E[Y_i(1)-Y_i(0)\mid \text{complier}].
\]
The denominator is
\[
    \E[D_i\mid Z_i=1]-\E[D_i\mid Z_i=0]
    =
    \E[D_i(1)-D_i(0)]
    =
    \Pr(\text{complier}).
\]
Dividing gives the complier average treatment effect.
\end{proof}

\subsection{LATE is not automatically ATE}

If treatment effects are homogeneous, then LATE, ATE, ATT, and other averages
coincide. But if treatment effects vary across people, IV estimates the effect
for the people induced by the instrument. This may be exactly the policy group
of interest, or it may be a narrow group.

\begin{boxF}
\textbf{Interpretation warning.} Do not write ``IV estimates the treatment
effect'' without saying for whom. In a non-compliance setting, the honest
interpretation is usually ``the effect for compliers induced by this
instrument.''
\end{boxF}

\subsection{Interpreting compliers}

Compliers are not observed one by one, but we can describe them indirectly. If
the instrument is distance to a training center, compliers may be people whose
participation decision is sensitive to travel costs. If the instrument is a
lottery offer, compliers are people who attend when offered but do not attend
without the offer. Different instruments define different complier populations.

\begin{exmp}[Non-compliance numeric example]
Suppose $60\%$ of those offered a program attend and $20\%$ of those not
offered attend. The first stage is $0.40$, so under monotonicity the complier
share is $40\%$. If assignment raises average earnings by $2$ units, then
\[
    \LATE = \frac{2}{0.40}=5.
\]
The offer has an average effect of $2$ on the assigned population, but the
program's effect on the marginal people induced to attend is $5$.
\end{exmp}

\section{Reporting IV Results}

A good IV table and discussion should report more than the final 2SLS
coefficient.
\begin{itemize}
  \item Define the outcome, treatment, instrument, sample, and controls.
  \item Report the first stage, including the sign, magnitude, and
        first-stage $F$-statistic.
  \item Report the reduced form when it is substantively meaningful.
  \item Report the IV/2SLS estimate with appropriate standard errors.
  \item Explain the identifying assumptions in words.
  \item Discuss possible exclusion violations.
  \item State the population affected by the instrument, especially the
        complier group in LATE settings.
  \item Compare OLS, ITT, and IV carefully; do not treat disagreement as
        automatic evidence that IV is better.
\end{itemize}

\begin{boxD}
\textbf{Connection to the code.} The lecture script prints three blocks:
endogeneity and IV recovery, weak instruments, and non-compliance/LATE. In the
first block, OLS is biased because it omits $U$, while IV recovers the true
effect because $Z$ is randomly assigned. In the second block, the same true
effect becomes hard to estimate as the first stage weakens. In the third block,
the Wald and 2SLS estimates recover the true complier effect, not the population
ATE.
\end{boxD}

\section{Worked Problems}

\subsection{Problem 1: First stage, reduced form, and Wald}

A randomized encouragement increases program participation from $25\%$ in the
control group to $55\%$ in the encouraged group. Average outcomes are $12$ in
the control group and $15.6$ in the encouraged group.

\textbf{Solution.} The first stage is
\[
    0.55-0.25=0.30.
\]
The reduced form is
\[
    15.6-12=3.6.
\]
The Wald estimate is
\[
    \hat{\tau}_{Wald} = \frac{3.6}{0.30}=12.
\]
Under the IV/LATE assumptions, the treatment effect for compliers is $12$.

\subsection{Problem 2: Why treatment received is endogenous}

Suppose assignment to a scholarship is random, but actual college attendance is
not. Motivated students are more likely to attend even without a scholarship.
Explain why comparing attendees to non-attendees is not a randomized comparison.

\textbf{Solution.} Randomization applies to the offer $Z$, not necessarily to
attendance $D$. Attendance depends on motivation, costs, information, and
preferences. These variables can also affect earnings. Therefore
\[
    \E[\Yzero \mid D=1] \neq \E[\Yzero \mid D=0],
\]
so the observed difference by treatment received mixes the treatment effect
with selection.

\subsection{Problem 3: Deriving the IV estimand}

Assume
\[
    Y_i=\alpha+\tau D_i+U_i,
    \qquad \Covp(Z_i,U_i)=0,
    \qquad \Covp(Z_i,D_i)\neq 0.
\]
Show that $\tau=\Covp(Z_i,Y_i)/\Covp(Z_i,D_i)$.

\textbf{Solution.} Taking covariance with $Z_i$ gives
\[
    \Covp(Z_i,Y_i)
    =
    \Covp(Z_i,\alpha+\tau D_i+U_i)
    =
    \tau\Covp(Z_i,D_i)+\Covp(Z_i,U_i).
\]
The constant has zero covariance with $Z_i$, and instrument exogeneity implies
$\Covp(Z_i,U_i)=0$. Divide by $\Covp(Z_i,D_i)$.

\subsection{Problem 4: Compliance shares}

In an experiment, $E[D\mid Z=1]=0.72$ and $E[D\mid Z=0]=0.18$. Assume
monotonicity. What share of the population are compliers?

\textbf{Solution.} The complier share is the first stage:
\[
    0.72-0.18=0.54.
\]
So $54\%$ of the population are compliers. This does not mean we know exactly
which individuals are compliers.

\subsection{Problem 5: Weak instrument intuition}

Two instruments are proposed for the same treatment. Instrument A changes
treatment probability by $0.40$. Instrument B changes treatment probability by
$0.02$. Both are argued to be random. Which one creates the bigger weak-IV
concern?

\textbf{Solution.} Instrument B. The IV estimator divides by the first stage.
With a first stage of only $0.02$, sampling noise in the reduced form or first
stage can produce very large changes in the ratio. Instrument B may still be
valid in principle, but it is weak and requires special caution.

\section{Checklist for Students}

Before trusting an IV estimate, answer these questions.
\begin{enumerate}
  \item What is the treatment $D$?
  \item What is the instrument $Z$?
  \item What outcome $Y$ is being studied?
  \item Why does $Z$ affect $D$? This is relevance.
  \item Is the first stage large enough to be useful?
  \item Why is $Z$ as good as randomly assigned, possibly conditional on $X$?
  \item Could $Z$ affect $Y$ through any channel other than $D$?
  \item Are controls pre-determined, or are any controls post-treatment?
  \item Is this a constant-effect IV setting or a LATE setting?
  \item If it is LATE, who are the likely compliers?
  \item Are standard errors computed using a proper IV routine?
  \item Does the paper report first stage, reduced form, and IV estimates?
\end{enumerate}

\section{Practice Questions}

\begin{enumerate}
  \item Explain the difference between an instrument and a control variable.
  Use an example where the same variable might be a bad instrument but a useful
  control.

  \item A judge's tendency to be strict is used as an instrument for sentence
  length. State the relevance, independence, and exclusion assumptions in
  words. Which assumption worries you most?

  \item In a scholarship experiment, the offer raises attendance from $50\%$
  to $80\%$ and raises earnings from $30$ to $33$. Compute the ITT on earnings,
  the first stage, and the LATE.

  \item Why is it wrong to interpret the LATE as the ATE when treatment effects
  are heterogeneous?

  \item Suppose the first-stage $F$-statistic is $3$. What does this suggest?
  What should the researcher report or do next?

  \item Give an example of an exclusion restriction violation for distance to
  college as an instrument for education.

  \item In a randomized encouragement design, define never-takers, compliers,
  always-takers, and defiers using $D(0)$ and $D(1)$.

  \item Run \texttt{lectures/code/lecture-5-iv-late.py}. In the weak-instrument
  block, what happens to the spread of IV estimates as the first-stage strength
  falls? Explain using the Wald ratio.
\end{enumerate}

\end{document}
